\documentclass[a4paper]{article}

\usepackage{graphicx}
\usepackage{amsmath,amssymb}
\usepackage{minted}
\usepackage{multirow}
\usepackage{float}
\usepackage{todonotes}
\usepackage{forest}
\usepackage{adjustbox}
\usepackage{xstring}
\usepackage{xcolor}
\usepackage[most]{tcolorbox}
\usepackage{xparse}
\usepackage{natbib}

\usetikzlibrary{graphs,graphdrawing}
\usegdlibrary{layered}

% For external Graphviz
\input{/home/donkarlo/Dropbox/repo/nd_ai_project/src/nd_ai/knoweldge_representation/language/formal/computer/programming/tex/latex/code_clips/external_graphviz.tex}

% For tags
\input{/home/donkarlo/Dropbox/repo/nd_ai_project/src/nd_ai/knoweldge_representation/language/formal/computer/programming/tex/latex/parts/control_sequence/kind/semtag/semtag.tex}

% For links
\input{/home/donkarlo/Dropbox/repo/nd_ai_project/src/nd_ai/knoweldge_representation/language/formal/computer/programming/tex/latex/code_clips/seehere.tex}

\title{Introduction to Kubernetes}
\date{}
\author{Mohammad Rahmani}

\begin{document}
    
    \maketitle
    
    
    \section{What Is Kubernetes?}
        
        Kubernetes, commonly abbreviated as \textbf{K8s}, is an open-source platform
        for automatically deploying, managing, and scaling containerized applications.
        
        A container packages an application together with its dependencies. Docker is
        commonly used to create and run such containers. However, when an application
        contains many containers or runs on multiple servers, manually managing the
        containers becomes difficult.
        
        Kubernetes solves this problem by coordinating the containers and deciding:
        
        \begin{itemize}
            \item where each container should run;
            \item how many copies of each container are required;
            \item what should happen when a container crashes;
            \item how network traffic should be distributed;
            \item how a new version of an application should be deployed.
        \end{itemize}
        
        \begin{tcolorbox}[
            enhanced,
            colback=blue!5,
            colframe=blue!50!black,
            title=Simple Definition
        ]
            Docker creates and runs containers, while Kubernetes manages many containers
            across one or more machines.
        \end{tcolorbox}
    
    
    \section{Why Is Kubernetes Needed?}
        
        Consider an online store consisting of the following components:
        
        \begin{itemize}
            \item a frontend application;
            \item a backend API;
            \item a database;
            \item a Redis cache;
            \item several background workers.
        \end{itemize}
        
        Each component can run inside a separate Docker container. For a small system,
        these containers can be managed manually or with Docker Compose.
        
        For a large production system, however, several problems appear:
        
        \begin{itemize}
            \item A container may unexpectedly stop.
            \item A server may become unavailable.
            \item The number of users may suddenly increase.
            \item Several copies of the backend may be required.
            \item A new application version must be deployed without long downtime.
            \item Requests must be distributed between several backend containers.
        \end{itemize}
        
        Kubernetes automates these operations.
    
    
    \section{Example}
        
        Assume that a backend application initially runs inside one container:
        
        \begin{minted}[frame=single,fontsize=\small]{text}
Backend Container 1
        \end{minted}
        
        When the number of users increases, Kubernetes can create additional copies:
        
        \begin{minted}[frame=single,fontsize=\small]{text}
Backend Container 1
Backend Container 2
Backend Container 3
Backend Container 4
        \end{minted}
        
        Kubernetes then distributes incoming requests among these containers.
        
        If \texttt{Backend Container 2} crashes, Kubernetes detects the failure and
        creates a replacement container. This behavior is called
        \textbf{self-healing}.
    
    
    \section{Main Responsibilities}
        
        \subsection{Container Deployment}
            
            Kubernetes starts the required containers according to a configuration
            provided by the developer.
        
        \subsection{Self-Healing}
            
            If a container crashes or becomes unhealthy, Kubernetes restarts or replaces
            it automatically.
        
        \subsection{Scaling}
            
            Kubernetes can increase or decrease the number of application instances.
            
            For example:
            
            \[
                N_{\text{backend}} =
                \begin{cases}
                    2, & \text{during normal traffic},\\
                    10, & \text{during high traffic}.
                \end{cases}
            \]
            
            Scaling may be performed manually or automatically according to CPU usage,
            memory usage, or other metrics.
        
        \subsection{Load Balancing}
            
            When several copies of an application are running, Kubernetes distributes
            network requests among them.
        
        \subsection{Rolling Updates}
            
            Kubernetes can gradually replace an old application version with a new one.
            This makes it possible to update an application with little or no downtime.
        
        \subsection{Service Discovery}
            
            Containers may be created, destroyed, or moved between machines. Kubernetes
            provides stable service names so that application components can communicate
            without knowing the exact address of every container.
    
    
    \section{Important Kubernetes Concepts}
        
        \subsection{Cluster}
            
            A Kubernetes cluster is a collection of machines managed as one system.
            
            A cluster normally contains:
            
            \begin{itemize}
                \item a \textbf{control plane}, which manages the cluster;
                \item one or more \textbf{worker nodes}, which run applications.
            \end{itemize}
        
        \subsection{Node}
            
            A node is a physical or virtual machine belonging to the Kubernetes cluster.
        
        \subsection{Pod}
            
            A Pod is the smallest deployable unit in Kubernetes.
            
            A Pod usually contains one application container, although it may contain
            multiple closely related containers.
            
            \begin{tcolorbox}[
                enhanced,
                colback=yellow!8,
                colframe=yellow!50!black,
                title=Important
            ]
                Kubernetes does not directly manage individual containers. It normally manages
                Pods, and each Pod contains one or more containers.
            \end{tcolorbox}
        
        \subsection{Deployment}
            
            A Deployment describes how an application should run.
            
            It can specify:
            
            \begin{itemize}
                \item the container image;
                \item the required number of Pod replicas;
                \item the application version;
                \item the update strategy.
            \end{itemize}
        
        \subsection{Service}
            
            A Service provides a stable network address for a group of Pods.
            
            Even if the Pods are replaced or receive new IP addresses, the Service remains
            available through the same name and address.
    
    
    \section{Basic Architecture}
        
        A simplified Kubernetes architecture can be represented as follows:
        
        \begin{center}
            \begin{adjustbox}{max width=\textwidth}
                \begin{forest}
                    for tree={
                        draw,
                        rounded corners,
                        align=center,
                        edge={->},
                        l sep=12mm,
                        s sep=6mm,
                        inner xsep=5pt,
                        inner ysep=4pt,
                        font=\small
                    }
                    [Kubernetes Cluster
                    [Control Plane
                    [API Server]
                    [Scheduler]
                    [Controller Manager]
                    ]
                    [Worker Node 1
                    [Pod
                    [Backend Container]
                    ]
                    [Pod
                    [Worker Container]
                    ]
                    ]
                    [Worker Node 2
                    [Pod
                    [Backend Container]
                    ]
                    [Pod
                    [Redis Container]
                    ]
                    ]
                    ]
                \end{forest}
            \end{adjustbox}
        \end{center}
        
        The control plane decides where applications should run. The worker nodes run
        the actual Pods and containers.
    
    
    \section{A Simple Kubernetes Configuration}
        
        The following configuration creates a Deployment containing three replicas of
        an Nginx web server:
        
        \begin{minted}[frame=single,fontsize=\small]{yaml}
apiVersion: apps/v1
kind: Deployment

metadata:
  name: nginx-deployment

spec:
  replicas: 3

  selector:
    matchLabels:
      app: nginx

  template:
    metadata:
      labels:
        app: nginx

    spec:
      containers:
        - name: nginx
          image: nginx:latest
          ports:
            - containerPort: 80
        \end{minted}
        
        The important line is:
        
        \begin{minted}[frame=single,fontsize=\small]{yaml}
replicas: 3
        \end{minted}
        
        It instructs Kubernetes to keep three Nginx Pods running.
        
        If one Pod crashes, Kubernetes creates another Pod so that the desired number
        remains three.
    
    
    \section{Docker and Kubernetes}
        
        Docker and Kubernetes are related, but they solve different problems.
        
        \begin{table}[H]
            \centering
            \begin{tabular}{
                |p{0.25\textwidth}
                |p{0.30\textwidth}
                |p{0.30\textwidth}|
            }
                \hline
                \textbf{Feature}
                &
                \textbf{Docker}
                &
                \textbf{Kubernetes}
                \\
                \hline
                
                Main purpose
                &
                Building and running containers
                &
                Managing containerized applications
                \\
                \hline
                
                Typical scale
                &
                One machine or a small environment
                &
                Multiple machines and many containers
                \\
                \hline
                
                Automatic recovery
                &
                Limited
                &
                Automatically replaces failed Pods
                \\
                \hline
                
                Automatic scaling
                &
                Not normally provided by Docker alone
                &
                Supports horizontal and automatic scaling
                \\
                \hline
                
                Load balancing
                &
                Requires additional configuration
                &
                Built-in through Services and related resources
                \\
                \hline
                
                Rolling updates
                &
                Usually managed manually
                &
                Built-in Deployment functionality
                \\
                \hline
            \end{tabular}
            
            \caption{A simplified comparison of Docker and Kubernetes.}
        \end{table}
    
    
    \section{Docker Compose and Kubernetes}
        
        Docker Compose is useful for running several related containers on one
        machine.
        
        For example, it can run:
        
        \begin{itemize}
            \item one backend container;
            \item one database container;
            \item one Redis container.
        \end{itemize}
        
        Kubernetes is more suitable when:
        
        \begin{itemize}
            \item the application runs on multiple servers;
            \item high availability is required;
            \item containers must restart automatically;
            \item automatic scaling is required;
            \item application updates must avoid downtime;
            \item the system contains many services.
        \end{itemize}
    
    
    \section{When Should Kubernetes Be Used?}
        
        Kubernetes is useful for:
        
        \begin{itemize}
            \item large web applications;
            \item microservice architectures;
            \item cloud-based services;
            \item applications with unpredictable traffic;
            \item systems requiring high availability;
            \item machine-learning services with several models and APIs;
            \item applications deployed across several servers.
        \end{itemize}
        
        Kubernetes may be unnecessary for:
        
        \begin{itemize}
            \item a small personal project;
            \item a simple website running on one server;
            \item an application with only one or two containers;
            \item a project for which Docker Compose already provides everything
            required.
        \end{itemize}
    
    
    \section{Summary}
        
        Kubernetes is a system for orchestrating containerized applications.
        
        Its main capabilities include:
        
        \begin{itemize}
            \item deploying containers;
            \item restarting failed applications;
            \item scaling applications;
            \item distributing network traffic;
            \item managing applications across multiple servers;
            \item updating applications with minimal downtime.
        \end{itemize}
        
        The central distinction is:
        
        \begin{tcolorbox}[
            enhanced,
            colback=green!5,
            colframe=green!45!black
        ]
            \centering
            
            \textbf{Docker packages and runs containers.}
            
            \medskip
            
            \textbf{Kubernetes coordinates and manages containers at scale.}
        \end{tcolorbox}

% \bibliographystyle{plainnat}
% \bibliography{/home/donkarlo/Dropbox/repo/nd_shared_research/bibliography/bibliography.bib}

\end{document}